\section*{A section without number}

You can achieve this with

% Set style for all upcoming listings
\lstset{style=mystyle} 

\begin{lstlisting}[language=TeX,%
  caption={Example of a section command without numbering},%
  label=lst:latex_section_no_number]
  \section*{A section without number}
\end{lstlisting}
  

\section{Figures \& Codes}
\label{sec:fig}

\begin{figure}[thb]
  \centering

  \stackunder{%
    \includegraphics[width=50mm,scale=0.5]{images/be_cool_and_learn_latex}}%
    {\scriptsize Source: \url{http://KeepCalmAndPosters.com/}%
  }

  \caption{Latex}
  \label{fig:learnLatex}
\end{figure}

From the picture above \ref{fig:learnLatex} we have decided to learn latex and
be COOL!

You can also use vref: \vref{fig:learnLatex}.

In order to add specific code in latex we have two ways of achieving it.
The following code is written in latex but interpreted as python code,


Here is some Python code:

\begin{lstlisting}[language=Python, caption= Python Example Code]
  num = float(input("Enter a number: "))
  if num > 0:
     print("Positive number")
  elif num == 0:
     print("Zero")
  else:
     print("Negative number")
\end{lstlisting}

In this example the package xcolor used in the \textbf{main.tex} is imported
and then the command \textbf{definecolor{}{}{}} is used to define new colours
in rgb format that will later be used. 

There are essentially two commands that generate the style for this example:

\textbf{lstdefinestyle{mystyle}{...}}

Defines a new code listing style called "\textit{mystyle}".
Inside the second pair of braces the options that define this style are passed.

The second method to display the python code without writing the code is by
importing the python file. 

\lstinputlisting[language=Python]{PrimeNumber.py}

To produce an inline expression, place the math expression between dollar signs
(\$).
For example: $E=mc^2$ or a mathematical expression
\[x=\frac{-b\pm\sqrt{b^2-4ac}}{2a}\]

If it is needed to express functions on a set we can proceed with the following
expression.

\begin{equation*} 
  h^1(s,g) = \begin{cases} 

    0, &  \text{if} g \subseteq s \\ 
    \min_{a \in a,t^\text{-1}(p_\text{new},a)} :=~ c(a) + h^1(s,t^\text{-1}(p_\text{new},a)), & |g| = 1 \\
    \max_{g^1 \in g} :=~ h^1 (s,\{g^1\}) & |g| > 1

  \end{cases} 
\end{equation*}
